\begin{chapintro}
\kkref{sec:13-5}で出した主解析は、カットオフ 0.5 ng/mL の 6 研究から感度 0.746、特異度 0.757 の要約点を出しました。それでも感度は研究によって 0.56 から 0.88 まで開いています。ばらつきの原因の候補を挙げたうえで、全 10 研究にカットオフを共変量として入れて調べ、バイアスのリスクが高い研究を除いた感度分析も行います。場面をカットオフで調整して比べる解析、検査どうしの比較、同じ PCT を扱った 2 本のメタ解析（Tang 2007、Wacker 2013）の読み比べは、発展の印をつけて載せました。
\end{chapintro}

\section{ばらつきの原因の候補}
\label{sec:14-1}

診断精度の研究が互いに違う結果を出す理由は、2×2 表の作られ方をさかのぼると整理できます。誰を対象にしたか、検査をどう測ったか、正解（参照基準）をどう決めたか、の 3 つです。

\begin{description}
\item[対象] 救急外来の発熱患者と ICU の重症患者では、菌血症の人も菌血症でない人も、病態が違います。ICU には菌血症でなくても手術や外傷で炎症の強い患者が多く、PCT が上がる理由が菌血症以外にもあります。重症度、年齢、腎機能、併存疾患も研究ごとに違います。
\item[対象検査] 測定試薬と機器、検体の種類（血清か血漿か）、採血の時点、カットオフが研究ごとに違います。
\item[参照基準] 血液培養を何セット採ったか、汚染菌と判断した陽性をどう扱ったか、抗菌薬投与後の培養を含めたかで、「菌血症あり」の人が入れ替わります。
\end{description}

対象検査の違いには、検査技師が最も具体的に気づける点があります。PCT の測定試薬は複数あり、同じ検体を測っても試薬によって値がそろわないことがあります。その場合、論文に同じ「0.5 ng/mL」と書かれていても、指している濃度が研究ごとに少しずつ違います。分析性能（測定値の精度や正確さ）の差が、臨床的な診断精度（感度と特異度）の差として現れることがあります\kksee{sec:09-6}。

測定試薬の違いを共変量にした例は、Cochrane の DTA ハンドブックにもあります。関節リウマチの診断に使う抗 CCP 抗体の測定試薬の世代の違いで、感度と特異度が違うかを二変量モデルで調べています（第 9 章と第 10 章）。

採血の時点も、診断精度に関わります。PCT は細菌感染が始まってから数時間かけて上昇するので、発熱した直後に採った検体では、菌血症の人でもまだ値が低いことがあります。来院時の 1 回だけを測った研究と、24 時間以内の最高値を使った研究では、感度が違うと予想されます。

挙げた候補は、データを見てから探すものではありません。\kkref{sec:02-5}で述べたとおり、どの特徴でばらつきを調べるかはプロトコルの段階で決め、データ抽出の表にその列を用意しておきます。データを見てから都合のよい共変量を探すと、偶然の差を原因と取り違えます。

\section{カットオフを共変量にする}
\label{sec:14-2}

二変量モデルに共変量を入れると、共変量の値ごとに要約点が動くモデルになります。介入のメタ解析のメタ回帰（短い説明は\kkref{sec:06-4}、係数の読み方は発展の\kkref{sec:06-5}）に当たる方法で、検査では感度と偽陽性割合の 2 つの値が同時に動きます。mada の reitsma() では、formula 引数に共変量を書きます。左辺の cbind(tsens, tfpr) は「logit 感度と logit 偽陽性割合の両方」を意味する決まった書き方です。

カットオフは連続した値なので、対数をとって共変量にします\kksee{sec:05-1-2}。0.25、0.5、1 ng/mL は、対数をとると等間隔になります（底を 2 にすると −2、−1、0。R の log() は自然対数で、−1.39、−0.69、0）。つまり、カットオフを 2 倍にするごとに同じだけ効果が動く、と置いたことになります。図 14.1 の横軸で 3 つのカットオフが等間隔に並ぶのは、このためです。この解析にはカットオフの混ざった全 10 研究を使います。目的は、カットオフを動かしたときに感度と偽陽性割合がどう動くかを調べることで、10 研究から 1 つの要約点を出すことではありません。カットオフの混ざった研究から要約の感度と特異度を出すことは、臨床で解釈できず、他の場面にも一般化できないとして、Cochrane の DTA ハンドブックが避けるよう求めています（Macaskill 2023）。

\begin{rcode}
pct$lcut <- log(pct$cutoff)
fit_cut <- reitsma(pct,
  formula = cbind(tsens, tfpr) ~ lcut)
summary(fit_cut)
\end{rcode}

\begin{rout}
                  Estimate Std. Error      z Pr(>|z|)
tsens.(Intercept)    1.174      0.366  3.210    0.001
tsens.lcut          -0.064      0.434 -0.147    0.883
tfpr.(Intercept)    -1.639      0.237 -6.925    0.000
tfpr.lcut           -0.774      0.280 -2.770    0.006
\end{rout}

tsens の行が logit 感度、tfpr の行が logit 偽陽性割合です。偽陽性割合（1 − 特異度）を使うのは、ROC 平面の横軸に合わせる mada の決まりです。切片は log(カットオフ) が 0 のとき、つまり基準にしたカットオフ 1 ng/mL のときの logit 値です。lcut の行は、log(カットオフ) に掛かる係数です。

tfpr の lcut の係数 −0.774 は、カットオフが 2 倍になると logit 偽陽性割合が −0.774 × log(2) ≈ −0.54 動く、という意味です。p = 0.006 で、カットオフを上げると偽陽性割合が下がることがはっきり示されました。logit 感度の係数は −0.064 でほぼ 0 に近く、p = 0.88 です。このデータでは、カットオフを動かして動くのは主に偽陽性割合のほうでした。

係数から、カットオフごとの感度と特異度を計算します。観察されたカットオフは 0.25、0.5、1 ng/mL の 3 通りなので、その範囲だけを求めます。

\begin{rcode}
cf <- coef(fit_cut)      # 係数の表（行は切片と lcut）
lc <- log(c(0.25, 0.5, 1))
sens <- plogis(cf["(Intercept)", "tsens"] +
               cf["lcut", "tsens"] * lc)
spec <- 1 - plogis(cf["(Intercept)", "tfpr"] +
                   cf["lcut", "tfpr"] * lc)
round(data.frame(cutoff = c(0.25, 0.5, 1), sens, spec), 2)
\end{rcode}

\begin{rout}
  cutoff sens spec
1   0.25 0.78 0.64
2   0.50 0.77 0.75
3   1.00 0.76 0.84
\end{rout}

カットオフを 0.25 から 1 に上げると、特異度は 0.64 から 0.84 まで上がり、感度は 0.78 から 0.76 へとほとんど動きません。閾値効果として予想される向きですが、感度の変化は小さく、偶然のばらつきと区別できません。この当てはめを図にしたものが図 14.1 です。

\kkfig[.9\linewidth]{x2_cutoff_trend.pdf}{図 14.1　カットオフと感度（左）、偽陽性割合（右）。円の大きさは研究の人数、線は log(カットオフ) を共変量にした二変量モデルの当てはめ}

\begin{rcode}
cx <- seq(log(0.2), log(1.3), length.out = 100)
n  <- with(pct, TP + FP + FN + TN)
par(mfrow = c(1, 2))

plot(pct$cutoff, pct$TP / (pct$TP + pct$FN), log = "x",
     ylim = c(0, 1), pch = 21, cex = 0.8 + 2 * sqrt(n / max(n)),
     xlab = "カットオフ（ng/mL、対数目盛）", ylab = "感度")
lines(exp(cx), plogis(cf["(Intercept)", "tsens"] +
                      cf["lcut", "tsens"] * cx))

plot(pct$cutoff, pct$FP / (pct$FP + pct$TN), log = "x",
     ylim = c(0, 1), pch = 21, cex = 0.8 + 2 * sqrt(n / max(n)),
     xlab = "カットオフ（ng/mL、対数目盛）", ylab = "偽陽性割合")
lines(exp(cx), plogis(cf["(Intercept)", "tfpr"] +
                      cf["lcut", "tfpr"] * cx))
\end{rcode}

線の式は、logit 感度 = 切片 + 係数 × log(カットオフ) です。cx は log(カットオフ) の値の並びで、plogis() で割合に戻してから、exp(cx) をカットオフの位置として線を引いています。

横軸は対数目盛のカットオフで、0.25、0.5、1 ng/mL が等間隔に並びます。右の偽陽性割合は、カットオフを上げるにつれて線が下がり、点の並びと合っています。左の感度の線はほぼ水平です。0.5 ng/mL の 6 研究だけでも感度は 0.56 から 0.88 まで開いていて、カットオフでは説明できない研究間の差が残ることが分かります。線は研究間のばらつきを含めた二変量モデルから出したもので、点の単純な回帰線ではありません。

場面で精度が違うかを見たいときは、まずカットオフをそろえた研究どうしで比べます。カットオフが場面と結びついていると、場面の差とカットオフの差が区別できないからです。この例で場面とカットオフを同時にモデルに入れる解析は、次の節に発展として載せました。

\section{\hatten 場面をカットオフで調整して比べる}
\label{sec:14-3}

救急外来の研究と ICU の研究では、対象の病態が違います。場面を共変量に加えるときは、カットオフも同時にモデルに入れます。カットオフが場面と結びついていると、場面だけを入れたモデルでは、場面の差とカットオフの差が区別できないからです。この例の場面とカットオフの組み合わせを表 14.1 に示します。

\begin{kkcenter}
\kkfit{\begin{tabular}{lccc}
\toprule
 & 0.25 & 0.5 & 1.0 \\
\midrule
救急外来 & 2 & 4 & 0 \\
ICU      & 0 & 2 & 2 \\
\bottomrule
\end{tabular}}
\kkcap{表 14.1　場面とカットオフ（ng/mL）ごとの研究数}
\end{kkcenter}

救急外来の 6 研究は 0.25 と 0.5 ng/mL、ICU の 4 研究は 0.5 と 1 ng/mL を使っています。0.5 ng/mL では両方の場面の研究がそろっているので、カットオフを調整したうえで場面を比べる足がかりがあります。

\begin{rcode}
fit_set <- reitsma(pct,
  formula = cbind(tsens, tfpr) ~ setting + lcut)
summary(fit_set)
\end{rcode}

\begin{rout}
                      Estimate Std. Error      z Pr(>|z|)
tsens.(Intercept)        1.225      0.277  4.427    0.000
tsens.setting救急外来   -0.994      0.403 -2.463    0.014
tsens.lcut              -0.800      0.427 -1.873    0.061
tfpr.(Intercept)        -1.623      0.251 -6.479    0.000
tfpr.setting救急外来    -0.202      0.346 -0.585    0.559
tfpr.lcut               -0.924      0.381 -2.423    0.015
\end{rout}

場面の列は文字なので、R は文字の並び順で先に来る値（この例では ICU）を基準にし、もう一方の値（救急外来）との差を推定します。出力は次のように読みます。

\begin{description}
\item[tsens.setting救急外来] 同じカットオフで比べると、救急外来の研究では logit 感度が 0.994 低いという推定で、p = 0.014 でした。
\item[tfpr.setting救急外来] 救急外来の研究で logit 偽陽性割合が 0.202 低いという推定ですが、p = 0.56 で、差ははっきりしません。
\item[tsens.lcut と tfpr.lcut] カットオフの係数です。偽陽性割合の −0.924（p = 0.015）は、カットオフを上げると偽陽性割合が下がることを示します。感度の −0.800 は p = 0.061 で、場面を調整すると感度もカットオフとともに下がる向きに出ました。カットオフだけを入れたモデルでは −0.064 だったので、場面を入れたことで推定が大きく変わっています。
\end{description}

カットオフが 0.5 ng/mL の研究について、係数から感度と特異度を計算します。救急外来の列の係数は、ICU との差を表します。

\begin{rcode}
cs <- coef(fit_set)
f <- function(par, e) cs["(Intercept)", par] +
  e * cs["setting救急外来", par] + cs["lcut", par] * log(0.5)
round(rbind(
  ICU = c(sens = plogis(f("tsens", 0)),
          spec = 1 - plogis(f("tfpr", 0))),
  救急外来 = c(sens = plogis(f("tsens", 1)),
              spec = 1 - plogis(f("tfpr", 1)))), 2)
\end{rcode}

\begin{rout}
         sens spec
ICU      0.86 0.73
救急外来 0.69 0.77
\end{rout}

ICU の感度は 0.86、特異度は 0.73、救急外来の感度は 0.69、特異度は 0.77 です。調整すると、場面による違いは偽陽性割合ではなく感度に現れました。カットオフを調整する前の見え方（救急外来で特異度が低い）とは異なります。

この結果から、救急外来では PCT の感度が低い、と結論することはできません。理由は 3 つあります。カットオフを調整したうえで場面を直接比べられるのは、両方の場面の研究がそろう 0.5 ng/mL の 6 研究（救急外来 4、ICU 2）だけで、共変量を 2 つ入れると推定が不安定になります。感度の係数は、場面を入れるかどうかで −0.064 から −0.800 まで動いており、モデルの選び方に結果が左右されています。ICU の高い感度には、\kkref{sec:14-5}で見る高リスクの研究（部分的検証の Qureshi 2018）も含まれています。

共変量の解析は、研究の間の比較です。介入のメタ解析のサブグループ解析（\kkref{sec:06-4}）と同じく、ランダム化されていない観察研究の比較なので、ほかの特徴との交絡を避けられません。記載されていない特徴（重症度の分布、測定試薬）も研究ごとに違っています。研究数が少ないときに共変量をいくつも試すと、偶然の差を拾う危険も増えます。結果は「差の原因の候補が見つかった」として報告し、原因と断定しない書き方にします。

\section{同じカットオフの研究だけでまとめる}
\label{sec:14-4}

本書は、\kkref{sec:13-2}のとおり主解析をカットオフ 0.5 ng/mL の 6 研究（Kato 2010、Lopez 2011、Mori 2013、Park 2016、Qureshi 2018、Rossi 2019）にしました。臨床でよく使われるカットオフで、研究数が最も多いためです。どのカットオフを主解析にするかは、プロトコルで先に決めておきます。カットオフが混ざったデータの扱いには 3 つの方針があり（\kkref{sec:10-6}、Macaskill 2023、9.3.1）、全 10 研究は、要約点ではなく SROC 曲線と共変量の解析に使いました。

\begin{description}
\item[同じカットオフの研究だけで統合する] 臨床で使うカットオフについて、要約点を出します。研究数が減るのが欠点で、本書の主解析はこの方法です。
\item[SROC 曲線を示す] HSROC モデル、または二変量モデルから導いた SROC 曲線で、カットオフを動かしたときの感度と特異度の関係を示します。個々のカットオフの値は分かりません。
\item[複数のカットオフを同時に扱う方法] Steinhauser らの方法などがあります。本書では使いません。
\end{description}

\begin{hattenbox}
複数のカットオフを同時に扱う方法は、1 つの研究が複数のカットオフで報告しているデータをすべて使い、カットオフごとの感度と特異度を推定する手法です。Steinhauser らの方法（Steinhauser 2016）と、Jones らがカットオフへの依存を定量化する方法（Jones 2019）があります。Cochrane の DTA ハンドブック第 10 章は、前者を R の diagmeta パッケージで実行する例を載せています（10.7.1）。検査値が連続量であることを前提にしたモデルで、本書では名前の紹介にとどめます。
\end{hattenbox}

6 研究の要約点は、感度 0.746（0.638〜0.831）、特異度 0.757（0.726〜0.786）でした。要約点はこの 6 研究のカットオフ 0.5 ng/mL に対応する値で、0.25 や 1 ng/mL の値ではありません。全 10 研究から引いた SROC 曲線は、カットオフを変えたときの感度と特異度の関係を示す図として使い、曲線上の 1 点を要約点として報告することはしません。

\section{バイアスのリスクで分ける}
\label{sec:14-5}

\kkref{sec:12-7}の QUADAS-2 の判定では、カットオフ 0.5 ng/mL の 6 研究のうち 2 研究が高リスクでした。Rossi 2019 は対象検査の領域で、データを見てからカットオフを決めていました。Qureshi 2018 はフローとタイミングの領域で、参照基準による検証が一部の患者にしか行われていませんでした（部分的検証）。高リスクの研究を除いても結果が変わらないかを確かめるのが、バイアスのリスクによる感度分析です（感度分析の考え方は\kkref{sec:07-6}、判定の使い方は\kkref{sec:12-8}で説明しました）。

判定の列は、データの q\_pat、q\_idx、q\_ref、q\_flow（患者選択、対象検査、参照基準、フローとタイミング）に入れてあります。4 つの領域のどれかが高の研究を除き、残った 4 研究に主解析と同じ、相関を 0 にしたモデルを当てはめます。to\_long() は\kkref{sec:13-3}で定義した関数です。

\begin{rcode}
pct5 <- subset(pct, cutoff == 0.5)        # 主解析の 6 研究
low <- subset(pct5, q_pat != "高" & q_idx != "高" &
                    q_ref != "高" & q_flow != "高")
fit_ind <- function(d)                    # 相関を 0 にしたモデル（第 13 章の g と同じ式）
  glmer(cbind(y, n - y) ~ 0 + sens + spec +
          (0 + sens | study) + (0 + spec | study),
        data = to_long(d), family = binomial)
g_low <- fit_ind(low)
\end{rcode}

最後に、感度と特異度を逆 logit で戻し、研究間の標準偏差を取り出します。

\begin{rcode}
b <- fixef(g_low)
se <- sqrt(diag(as.matrix(vcov(g_low))))
plogis(b["sens"] + c(0, -1.96, 1.96) * se["sens"])
plogis(b["spec"] + c(0, -1.96, 1.96) * se["spec"])
sapply(VarCorr(g_low), attr, "stddev")   # 研究間の標準偏差を取り出す
\end{rcode}

\begin{rout}
[1] 0.709 0.573 0.815
[1] 0.781 0.755 0.806
  study.sens study.1.spec
       0.488        0.000
\end{rout}

上の 2 行は、推定値、95\% 信頼区間の下限、上限の順です。残った 4 研究（Kato、Lopez、Mori、Park）の要約点は、感度 0.709（0.573〜0.815）、特異度 0.781（0.755〜0.806）でした。6 研究の要約点（感度 0.746、特異度 0.757）と比べると、感度は下がり、特異度は上がっています。2 つの要約点を ROC 平面に置いたものが図 14.2 です。

\kkfig[.7\linewidth]{x2_rob_sens.pdf}{図 14.2　カットオフ 0.5 ng/mL の 6 研究。三角が高リスクの研究（Qureshi 2018、Rossi 2019）。青の菱形と楕円は 6 研究の要約点と 95\% 信頼領域、赤の菱形と破線の楕円は高リスクを除いた 4 研究の要約点と 95\% 信頼領域}

動いた向きは、2 つの高リスクの理由と整合します。データを見てカットオフを決めると、そのデータでの精度が高く出ます。部分的検証では、検査が陽性の人ほど参照基準で確かめられやすいので、感度が高く、特異度が低く出ます。Qureshi 2018 は感度 0.86 で、6 研究のなかで Park 2016（0.88）に次いで高く、特異度は 0.71 で最も低い側にあります。部分的検証で予想される向きです。Rossi 2019 は感度 0.73、特異度 0.71 で、個々の点からは偏りの向きを読み取れません。

高リスクの研究を除くと感度は下がり、特異度は上がりました。感度が高く見えていた分が取れたという向きですが、確かな結論とは言えません。2 つの要約点の信頼区間は大きく重なっており、図 14.2 でも 2 つの楕円は重なっています。4 研究では推定が不安定です。特異度の研究間の標準偏差は 0.000 と推定され、モデルが境界に張り付きました（出力に「boundary (singular) fit」の警告が出ます）。4 研究の特異度のばらつきが二項分布の誤差だけで説明できる範囲と推定されたということで、本当にばらつきがないことを示すものではありません。特異度の信頼区間は、研究間のばらつきを含まない分だけ狭く出ていると考えます。

mada の reitsma() で同じ 4 研究を推定すると、感度 0.706（0.546〜0.827）、偽陽性割合 0.222（0.196〜0.251）で、感度の向きは同じです。報告では、主解析を 6 研究の結果とし、感度分析で高リスクの研究を除いた結果と、境界に張り付いたことを並べて書きます。

\section{\hatten 検査どうしを比べる}
\label{sec:14-6}

PCT と CRP のどちらが菌血症を見分けるのに優れているか、という問いもよく立てられます。比べ方は 2 通りあり、信頼できる程度が違います。

\begin{description}
\item[直接比較] 同じ患者で PCT と CRP の両方を測った研究だけを使う方法です。患者、参照基準、場面がそろっているので、差は検査の違いによるものと考えやすくなります。
\item[間接比較] PCT の研究と CRP の研究を別々に集め、それぞれの要約点を比べる方法です。研究の対象や参照基準が違うので、差が検査の違いによるのか研究の違いによるのかを区別できません。
\end{description}

Cochrane の DTA ハンドブックは、比較精度のレビューでは直接比較の研究を優先するよう述べています（9.4.7）。間接比較には、検査の違いと研究の違いが混ざる問題があるためです。比較精度の研究のバイアスのリスクは、QUADAS-2 に比較の観点を加えた QUADAS-C で評価します（Yang 2021）。

二変量モデルでは、検査の種類を共変量として入れて比べます。このとき、同じ患者で両方を測った研究は、1 つの研究から 2 組の 2×2 表が出ます。2 組は同じ患者から作られているので互いに独立ではなく、ふつうの二変量モデルはこの依存を無視します。また、検査ごとに研究間のばらつきの大きさが違うことも多いので、Cochrane の DTA ハンドブック第 10 章は、検査ごとに研究間分散を別々に推定するモデルを示しています。直接比較の研究が十分にあれば、それだけで解析した結果を主に示し、間接比較を含めた結果は補助として扱います。

\section{\hatten 2 本の Lancet Infect Dis 論文}
\label{sec:14-7}

PCT の診断精度については、同じ Lancet Infectious Diseases 誌に、結論の違う 2 本のメタ解析が載っています。どちらも、重症患者で敗血症と非感染性の全身性炎症反応症候群（SIRS）を見分けられるか、という問いを扱いました。本書の例題（菌血症）とは標的疾患が違いますが、同じ検査について要約値と結論が変わった実例です。

\subsection{\hatten Tang 2007}
\label{sec:14-7-1}

Tang らは ICU の重症患者を対象にした 18 研究を統合し、感度と特異度がともに 71\%（95\% 信頼区間 67〜76\%）、SROC の AUC 0.78（0.73〜0.83）と報告しました。研究を Sackett と Haynes の分類で phase 2（14 研究）と phase 3（4 研究）に分けて調べ、phase 3 の研究では標本の大きさで結果が違い、小さい研究ほど精度が高く、大きい研究では差がない方向に近づいたとしています（メタ回帰の係数 −0.592、p = 0.017）。結論は「PCT は敗血症と非感染性 SIRS を確実には区別できない」というものでした。同誌の同じ年には、この論文への批判の投書が 3 通と、著者の返答が掲載されています。

\subsection{\hatten Wacker 2013}
\label{sec:14-7-2}

Wacker らは 7 つのデータベースを 2012 年 2 月まで検索し、3487 件から 30 研究、3244 人を組み入れました。二変量モデルで統合した感度は 0.77（0.72〜0.81）、特異度は 0.79（0.74〜0.84）、AUC は 0.85（0.81〜0.88）でした。I² は 96\% で、対象、入院の区分、測定法、重症度、参照基準の記述と盲検化のいずれのサブグループでも、ばらつきを説明できませんでした。結論は「PCT は早期診断に役立つが、病歴、診察、微生物学的な評価と合わせて解釈する」というものです。

\subsection{\hatten 読み比べて分かること}
\label{sec:14-7-3}

2 本の要約値の差は、感度で 71\% 対 77\%、特異度で 71\% 対 79\%、AUC で 0.78 対 0.85 でした。2 本の間で、組み入れた研究の数（18 と 30）、検索した時期、組み入れの基準が違います。その違いのどれが結論を分けたのかは、抄録から読み取れる範囲では決められません。

どちらのレビューでも研究間のばらつきが大きく、Tang は小さい研究ほど精度が高く出る傾向を示し、Wacker では測定法を含むどの特徴でもばらつきを説明できませんでした。ばらつきが大きく原因が分からないときは、要約点の数値の差を比べるより、新しい現場で精度がどれだけ振れるか（予測領域の広さ）を見るほうが臨床の判断に役立ちます\kksee{sec:13-6}。

\section{既存のメタ解析を読むとき}
\label{sec:14-8}

既存のメタ解析を読むときは、要約点がどのカットオフの値なのかを、方法の節で確かめます。研究ごとにカットオフが違うのに 1 つの要約点を出していれば、本書が\kkref{sec:13-2}で避けた読み方になります。たとえば Wacker 2013 は 30 研究を二変量モデルで 1 組の感度と特異度にまとめていますが、抄録にはその要約点がどのカットオフに当たるかが書かれていません。

自分でレビューを行うときは、標的疾患と対照群（菌血症か敗血症か、比べる相手は誰か）、組み入れる研究デザイン、カットオフの扱い、ばらつきを調べる共変量（測定法、場面、採血の時点）、小規模研究効果の調べ方を、プロトコルで先に決めておきます。それらが書いてあれば、結論が先行研究と違ったときに、どこで違ったのかを読者が追えます。

\section{この章のまとめ}
\label{sec:14-9}

研究間のばらつきの原因は、対象、対象検査、参照基準の 3 つの方向から候補を挙げ、プロトコルで調べる共変量を決めておきます。共変量は二変量モデルに入れ、感度と偽陽性割合のそれぞれへの影響を読みます。PCT の全 10 研究では、カットオフを上げると偽陽性割合が下がりました。場面を比べるときは、まずカットオフをそろえた研究どうしで比べます。カットオフの混ざった研究から要約点は出さず、主解析は同じカットオフの研究だけで行います。高リスクの研究を除いた感度分析で、感度は下がり特異度は上がりましたが、区間は重なっており、4 研究の推定は不安定でした。既存のメタ解析を読むときは、要約点がどのカットオフの値かを確かめます。\kkchref{15}では、主解析の要約点を 1000 人あたりの数に直し、確実性を評価して報告します。
